\chapter{Introduction}
\label{ch:introduction}

A flying robot is a body and a controller. The body decides what the robot can do: how much lift its wing produces, how much energy it spends per metre, how quickly it can turn. The controller decides what the robot does with that body. The two choices depend on each other. The best controller for one body is not the best controller for another body, and the value of a body can only be evaluated through the flight that a controller makes of it. Despite this coupling, flying robots are usually designed in sequence. The body comes first, from an engineer or from a commercial product, and the controller is designed afterwards for that body. This thesis studies the alternative, in which the body and the controller are designed together. It does so for a winged drone that has to fly through a forest, and it proposes a controller that adapts to the body it is given while it flies.

The drone of this thesis is a fixed-wing glider with a propeller at its nose and a wing in two panels. Each panel sweeps backwards and twists about its axis under a servomotor, and the elevator and the rudder move as a whole. The drone has no ailerons: it steers by changing the shape of its wing. Its body is described by fifteen parameters: the span, the aspect ratio and the airfoil section of the wing, the length of the fuselage, the position of the centre of mass and of the wing along the fuselage, the size of the elevator and of the rudder, the dihedral, and the gear ratios of the sweep and twist joints. These parameters span a wide space of bodies, from light and agile drones of about half a kilogram to gliders of two kilograms with a wing of almost a square metre. The controller is a neural network. It reads the altitude, the attitude and the velocity of the drone and a coarse depth image of the trees ahead, and it sets the throttle and the six joints twenty-five times per second. The task is to fly as far as possible through a corridor of trees whose density grows along the way, and to spend as little energy as possible in doing so. A flight is evaluated by two quantities. The \emph{progress} is the distance covered along the corridor. The \emph{cost of transport} is the energy spent per unit of weight and of distance. Every flight of the thesis takes place in simulation, in a physics engine that advances tens of thousands of drones in parallel on a graphics card.

\section{Co-Design and Its Cost}
\label{sec:intro-codesign}

Designing the body and the controller together is called \emph{co-design}. As an optimization problem, co-design is nested. An outer search runs over the bodies, and it scores every candidate body by the flight of the best controller for that body. An inner search finds that controller. The inner search is what makes the problem expensive. The co-design works that use learned controllers train a new controller from scratch for every body the search visits \cite{gupta2021embodied, muff2026hexacopters}. In the setting of this thesis, a controller trained on many bodies takes fourteen hours on a computing node with two graphics cards. A controller trained on a single body takes less than two hours on one of the cards. One search over bodies evaluates 1\,600 distinct bodies. A training for each of those bodies would take two and a half years in the first case. In the second case it would take almost two months, with two trainings running at a time. The works that take this route stay under a thousand bodies in total, or they require more than a thousand processors for a single run.

The cheap alternative is to train one \emph{generalist} controller on many bodies at once, and to let that single controller score every candidate body. This alternative has a cost of another kind. The generalist is a compromise between all the bodies it was trained on. It flies well the bodies that resemble its training distribution, and it flies poorly the bodies that differ from that distribution. The bodies that differ are the bodies that a search should find, because a body better than the average body must differ from the average body in size, in shape or in control authority. Two bodies can therefore come in one order under the generalist and in the opposite order under their own controllers. A search under the generalist then keeps the body that suits the generalist, and it discards the body with the greater potential. Training a controller for the chosen body afterwards does not repair this choice, because the body that was discarded cannot be recovered.

These facts are presented in \autoref{ch:challenges} as three challenges. The sequential design challenge: a body chosen before any search can never be improved by the search, whatever controller is trained on it. The morphology optimization challenge: a search under one fixed controller finds the bodies that suit that controller, not the bodies with the greatest potential. The computational feasibility challenge: training a controller for every body the search visits costs more than any search can afford. Together, the three challenges fix the shape of any workable solution. \emph{The controller must be trained once, for the whole space of bodies, and its adaptation to the body being scored must come from a mechanism that costs no more than the flights that score that body.} Whether an adaptation that cheap can change which bodies the search finds is the empirical question of this thesis.

\section{Plasticity as Adaptation}
\label{sec:intro-plasticity}

The answer proposed in this thesis is a controller whose weights change during flight under evolved Hebbian rules. The starting point is a generalist: a recurrent neural network, with a Long Short-Term Memory layer followed by two feed-forward layers and a linear output layer, trained by reinforcement learning with Proximal Policy Optimization on 256 random bodies at once. The controller is never told which body it is flying. It receives no description of the body, so that it learns to fly any morphology without hints on what the flight dynamics could be. Once trained, the network is frozen. Only the weights of its output layer are allowed to change: 224 gains that translate the 32 features of the last hidden layer into the seven commands. At every control step each gain moves by an amount that depends on the activity of the feature it reads and of the command it drives. The dependence is the generalized form of Hebb's postulate, with four coefficients for each gain. Each gain also decays back towards its trained value. The 896 coefficients are not learned by gradient descent. They are evolved with the Covariance Matrix Adaptation Evolution Strategy (CMA-ES), on the mean reward of the flights of a whole population of bodies. The rules are shown neither the body nor the reward. The body nonetheless reaches the rules, because the body is inside the control loop. The commands act on the body, the body answers according to its wing, its tail and its mass, and the frozen layers turn that answer into the activity that the rules read. Two bodies flown with the same rules therefore end the flight with different weights.

This use of plasticity differs from the one found in the literature on evolved plastic networks. There, the weights are drawn at random at the beginning of every episode, and the rules have to build a working controller from nothing \cite{floreano2000evolutionary, najarro2020meta, vandiggelen2025emergent}. In other works the fixed weights and the rules are co-trained by the same gradient \cite{miconi2018differentiable}. Here the controller already exists and already flies. \emph{Plasticity is not used to learn the task, but to adapt to the body a solution of the task that was learned for all bodies.} The rules start every flight from the trained weights and move them by a bounded amount. When every coefficient is zero, the Hebbian controller is the frozen generalist. This property gives the search of the rules a natural starting point, and it gives the experiments an exact control condition.

The Hebbian controller is placed inside a co-design pipeline made of two nested evolutionary searches. The inner loop evolves the rules with CMA-ES. At every generation, 64 candidate sets of rules fly the 64 bodies of the current population on 60 forests each, about a quarter of a million flights. The mean reward of its flights is the fitness of each set of rules. The outer loop evolves the bodies with the Non-dominated Sorting Genetic Algorithm II (NSGA-II). Every four generations of the inner loop, the eight best sets of rules fly every body on a longer course, the exam. The progress and the cost of transport measured in the exam are the two objectives of NSGA-II, which keeps half of the bodies and replaces the other half with offspring. The search of the rules is not restarted when the bodies change. A run of 200 generations examines 3\,136 bodies, of which at most 1\,600 are distinct, and takes about four days on one computing node, and no network is trained inside the run. The control condition is the same pipeline with the learning rate of the rules set to zero. The control starts from the same bodies, uses the same operators, the same exam and the same admission gate, and flies every body with the frozen generalist. The two conditions differ in the 896 coefficients, and the control is given the larger budget in computing time and in examined bodies.

\section{Contributions}
\label{sec:intro-contributions}

The contributions of this thesis are the following.
\begin{enumerate}
\item \textbf{A plastic controller on top of a frozen policy.} Evolved Hebbian rules act on the output layer of a policy that was pretrained by reinforcement learning and then frozen (\autoref{sec:hebbian-controller}). To the knowledge of the author, no previous work places evolved rules on top of a frozen pretrained policy, and no previous work uses plasticity across a population of bodies inside a co-design loop.
\item \textbf{A co-design pipeline that scores bodies by flights alone.} An evolution strategy over the rules is nested in time inside a multi-objective search over the bodies, and no controller is trained inside the search (\autoref{sec:pipeline}). The pipeline evaluates 1\,600 distinct bodies in four days on one computing node. A controller training per body would take months on the same node.
\item \textbf{A simulation setting for the co-design of winged drones.} The task, the procedurally generated forests, the parametric space of bodies and the measurement of the two objectives were set up on the Genesis engine \cite{genesis2024} and on an aerodynamic solver developed in the Laboratory of Intelligent Systems of EPFL, where this work was carried out (\autoref{ch:environment}). The setting flies a quarter of a million drones per generation on two graphics cards.
\item \textbf{A controlled comparison with morphology-only evolution.} Eight runs of the pipeline are compared with eight runs of the search under the fixed generalist, six of them started from the same bodies, with the larger budget given to the control (\autoref{sec:codesign-experiments}). The co-design fronts are up to 17\,\% cheaper at equal progress and fly up to 12\,\% farther at equal cost of transport. On these indicators every co-design run is better than every control run. The evolved bodies are different bodies: for three of the six seeds that carry a run of each kind the two runs end farther apart than control runs do among themselves, and the co-design runs end farther from the control runs than seven of the eight control runs end from one another.
\item \textbf{An analysis of what the rules do.} On a fixed body the rules move any base controller, the generalist and the specialists trained for that body alike, towards a slower and cheaper flight of the same course (\autoref{sec:inner-loop-experiments}). Inside the search the rules specialize to the family of bodies that the outer loop retains. They fly those bodies farther than the frozen generalist, they do not transfer that gain to bodies selected under the generalist, and they lose their generality along the run (\autoref{subsec:results-swapped-fronts} to \autoref{subsec:results-generality}).
\end{enumerate}

\section{Structure of the Thesis}
\label{sec:intro-structure}

The thesis is organized as follows. The technical background is given in \autoref{ch:background}: reinforcement learning and Proximal Policy Optimization, recurrent networks and the Long Short-Term Memory, evolutionary algorithms with CMA-ES and NSGA-II, and Hebbian learning with the plasticity rules that evolution can search. \autoref{ch:challenges} opens with the related work on co-design and on evolved plasticity, and with the position of this thesis with respect to both. It then states the three challenges, and the requirement they impose on the method. The simulation is described in \autoref{ch:environment}: the engine, the forest and the flight that is evaluated in it, and the drone with its space of bodies and its reference body. The method is described in \autoref{ch:method}: the base controllers and their training, the Hebbian controller with the reasons for making the output layer plastic, and the co-design pipeline with its control condition. The experiments are reported in \autoref{ch:results}: first the inner loop on fixed bodies, with the generalist and with the specialists as base, then the full pipeline against morphology-only evolution, with the analysis of the fronts, of the evolved bodies and of the evolved rules. The conclusions and three directions for future work are in \autoref{ch:conclusions}. The internal organization of the engine, the complete aerodynamic model and the derivations that the main text leaves out are collected in \autoref{app:simulation}.
